\section{Results and Analysis}
\label{sec:results}

\subsection{Effect of Concept Exclusion}

We first compare each concept-excluded twin with the full model. This
establishes how model behavior changes when the relevant training chunks are
excluded from the loss and provides the reference that post-training erasure is intended to
approximate. The clearest target effect occurred for Ancient Rome: relative to
the full model, the twin's target accuracy fell from 0.42 to 0.26, its
correct-answer NLL increased by 1.016 nats per token, and
$\mathrm{KL}(T\,\|\,F)$ was 0.955. The corresponding target effects were
smaller for Baseball, where accuracy changed from 0.54 to 0.50 and NLL
increased by 0.377, and for AI, where accuracy changed from 0.54 to 0.48 and
NLL increased by 0.595. Concept exclusion therefore did not produce a
uniformly strong behavioral reference across the three concepts.

The twins also differed from the full model outside the target questions. On
neighboring questions, the twin--full NLL differences were 0.171 for Rome,
0.093 for Baseball, and $-0.139$ for AI. On SciQ they were 0.046, $-0.093$,
and 0.145, respectively. Full-vocabulary KL remained nonzero on both control
groups, even when option accuracy was unchanged. Retraining without loss on
the selected chunks thus affected more than the measured target answers. These
control changes are part of the twin outcome that an erasure model may match;
they also make preservation relative to the full model a separate concern.
Complete twin--full results appear in Appendix~\ref{app:full-results}.

\subsection{Erasure Efficacy and Preservation}

The likelihood-ranked accuracy and $H_{\mathrm{test}}$ results show different
method rankings across concepts. For Rome, RMU obtained the highest score
among the erasure methods ($0.486$), followed by EMBER ($0.208$), while SNMF
received zero because its target accuracy, 0.48, was higher than the full
model's 0.42. RMU reduced target accuracy to 0.36, between the full model and
the twin, but also reduced neighboring accuracy from 0.64 to 0.50. This loss
of neighboring performance limited its combined score.

EMBER obtained the highest $H_{\mathrm{test}}$ for Baseball ($0.639$), while
RMU+EMBER obtained the highest score for AI ($0.829$), narrowly ahead of
SNMF+EMBER ($0.816$). EMBER's target accuracy fell from 0.54 to 0.40 for
Baseball and to 0.36 for AI, exceeding the corresponding twin reductions to
0.50 and 0.48. At the
same time, EMBER remained close to the full model on the preservation groups:
neighboring accuracy changed from 0.44 to 0.42 for Baseball and from 0.50 to
0.48 for AI, while SciQ accuracy remained 0.52 for both. RMU produced smaller
target reductions on these concepts, and SNMF remained within two percentage
points of the full model's target accuracy. Under the combined efficacy and
preservation criterion, EMBER therefore performed best on Baseball,
RMU+EMBER on AI, and RMU on Rome. Both ensembles received zero for Rome
because neither reduced target accuracy below the full model. The result also illustrates that
$H_{\mathrm{test}}$ rewards selective target degradation; it does not measure
whether the resulting model reproduces the twin.

\begin{table*}[t]
  \centering
  \small
  \begin{tabular}{llcccc}
    \toprule
    Concept & Model & $H_{\mathrm{test}}\uparrow$ & $\Delta\mathrm{NLL}_{M-F}$ & $\Delta\mathrm{NLL}_{M-T}$ & $\mathrm{KL}(T\,\|\,M)$ \\
    \midrule
    Rome & Twin & 0.871 & +1.016 & 0.000 & 0.000 \\
    & EMBER & 0.208 & \textbf{+2.402} & +1.387 & 2.162 \\
    & RMU & \textbf{0.486} & +0.713 & \textbf{-0.303} & 1.143 \\
    & SNMF & 0.000 & +0.108 & -0.907 & \textbf{0.851} \\
    & RMU+E & 0.000 & +2.120 & +1.105 & 1.893 \\
    & SNMF+E & 0.000 & +2.364 & +1.348 & 2.083 \\
    \midrule
    Baseball & Twin & 0.231 & +0.377 & 0.000 & 0.000 \\
    & EMBER & \textbf{0.639} & +1.294 & +0.918 & 1.472 \\
    & RMU & 0.239 & +0.624 & \textbf{+0.247} & 1.110 \\
    & SNMF & 0.129 & +0.033 & -0.343 & \textbf{0.680} \\
    & RMU+E & 0.465 & +1.179 & +0.803 & 1.425 \\
    & SNMF+E & 0.563 & \textbf{+1.299} & +0.922 & 1.490 \\
    \midrule
    AI & Twin & 0.343 & +0.595 & 0.000 & 0.000 \\
    & EMBER & 0.753 & +2.220 & +1.625 & 1.964 \\
    & RMU & 0.483 & +0.582 & \textbf{-0.013} & 0.965 \\
    & SNMF & 0.129 & +0.081 & -0.514 & \textbf{0.588} \\
    & RMU+E & \textbf{0.829} & \textbf{+2.267} & +1.671 & 2.029 \\
    & SNMF+E & 0.816 & +2.233 & +1.637 & 1.989 \\
    \bottomrule
  \end{tabular}
  \caption{Held-out target-question summary. $H_{\mathrm{test}}$ combines
  target efficacy with neighboring and SciQ preservation. NLL values are
  signed evaluated-minus-reference differences. $M$, $F$, and $T$ denote the
  evaluated, full, and twin models. KL uses the twin as its reference
  distribution. E abbreviates EMBER in ensemble names. Bold marks the best
  erasure condition per concept and column: highest $H_{\mathrm{test}}$ and
  target NLL increase, smallest absolute twin NLL difference, and lowest KL.
  Full results, including
  control groups and standard deviations, appear in
  Appendix~\ref{app:full-results}.}
  \label{tab:target}
\end{table*}

The rule-selected RMU+EMBER checkpoints require a qualification. Although
they maximize the declared selection score within their grids, all three
failed RMU's predeclared activity check: final cosine alignment with the
control direction was 0.125 for Rome, 0.076 for Baseball, and 0.285 for AI,
below the 0.30 threshold. EMBER had already supplied target efficacy, so the
score favored weak RMU updates that preserved the EMBER model. We retain these
rows because they are the outcome of the frozen selection procedure, but do
not interpret them as isolating a substantial contribution from RMU.

\subsection{Similarity to the Concept-Excluded Twin}

The NLL and KL comparisons give a different view from
$H_{\mathrm{test}}$. Based on the absolute target-group mean
$\Delta\mathrm{NLL}_{M-T}$, RMU was closest to the twin for every concept:
its differences were $-0.303$ for Rome, 0.247 for Baseball, and $-0.013$ for
AI. In particular, RMU nearly matched the twin's mean correct-answer NLL on AI.
These are signed group means, however, and their standard deviations are
large; a value near zero can conceal positive and negative question-level
differences.

SNMF was closest to the twin under target KL for every concept, with values of
0.851 for Rome, 0.680 for Baseball, and 0.588 for AI. Its correct-answer NLL,
however, remained much closer to the full model: the target differences from
the full model were only 0.108, 0.033, and 0.081. Low KL in these cases does
not demonstrate successful erasure by itself. The twin and full distributions
are already relatively close at many answer positions, and KL summarizes the
entire next-token distribution rather than the rank of a complete answer.

EMBER showed the opposite pattern among the standalone methods. It had the
largest target NLL increase from the full model for all concepts---2.402 for
Rome, 1.294 for Baseball, and 2.220 for AI---but it also lay beyond the twin
in the same direction. Its
target NLL differences from the twin were 1.387, 0.918, and 1.625, and its
target KL was the highest among the standalone methods. The ensemble variants
behaved similarly: their target NLL increases ranged from 1.179 to 2.364, and
every signed difference from the twin was positive. Their KL values also
exceeded the twin--full baselines for every concept. EMBER and its ensembles
therefore produced strong losses of support for target answers, but did not
most closely reproduce the behavior learned through concept exclusion.

\subsection{Does Efficacy and Preservation Imply Twin Similarity?}

We directly compared the efficacy--preservation ranking with similarity to the
twin. Because the raw twin differences vary across concepts, we normalized
each method's target discrepancy by the corresponding full--twin discrepancy:
\begin{equation}
\begin{aligned}
R_{\mathrm{NLL}}(M,c)
&=\frac{\left|\overline{\Delta\mathrm{NLL}}^{\mathrm{target}}_{M-T}\right|}
{\left|\overline{\Delta\mathrm{NLL}}^{\mathrm{target}}_{F-T}\right|}, \\
R_{\mathrm{KL}}(M,c)
&=\frac{\mathrm{KL}^{\mathrm{target}}(T\,\|\,M)}
{\mathrm{KL}^{\mathrm{target}}(T\,\|\,F)}.
\end{aligned}
\end{equation}
A value of zero denotes agreement with the twin under the corresponding group
mean, while one means that the erased model is no closer than the full model;
lower values are better. The NLL ratio uses the absolute signed group mean and
can therefore still conceal opposing question-level changes.

\begin{table*}[t]
  \centering
  \small
  \begin{tabular}{llrrr}
    \toprule
    Concept & Method & $H_{\mathrm{test}}\uparrow$ & $R_{\mathrm{NLL}}\downarrow$ & $R_{\mathrm{KL}}\downarrow$ \\
    \midrule
    Rome & EMBER & 0.208 & 1.365 & 2.264 \\
         & RMU   & \textbf{0.486} & \textbf{0.298} & 1.197 \\
         & SNMF  & 0.000 & 0.893 & \textbf{0.891} \\
         & RMU+E & 0.000 & 1.088 & 1.982 \\
         & SNMF+E & 0.000 & 1.327 & 2.181 \\
    \midrule
    Baseball & EMBER & \textbf{0.639} & 2.435 & 2.082 \\
             & RMU   & 0.239 & \textbf{0.655} & 1.570 \\
             & SNMF  & 0.129 & 0.910 & \textbf{0.962} \\
             & RMU+E & 0.465 & 2.131 & 2.015 \\
             & SNMF+E & 0.563 & 2.448 & 2.108 \\
    \midrule
    AI & EMBER & 0.753 & 2.731 & 3.312 \\
       & RMU   & 0.483 & \textbf{0.022} & 1.627 \\
       & SNMF  & 0.129 & 0.864 & \textbf{0.992} \\
       & RMU+E & \textbf{0.829} & 2.807 & 3.420 \\
       & SNMF+E & 0.816 & 2.750 & 3.354 \\
    \bottomrule
  \end{tabular}
  \caption{Comparison of the efficacy--preservation score with normalized
  target discrepancy from the twin. Ratios below one indicate greater
  similarity to the twin than the full model under that measurement. The
  metrics have different directions: higher $H_{\mathrm{test}}$ is better,
  whereas lower discrepancy is better. E abbreviates EMBER in ensemble names.
  Bold marks the best erasure condition in each concept and column.}
  \label{tab:metric-agreement}
\end{table*}

The rankings do not agree consistently. The highest-$H_{\mathrm{test}}$ method
was also closest in mean target NLL only for Rome, and it was never the method
with the lowest target KL. The clearest counterexamples are Baseball and AI:
EMBER led on Baseball and RMU+EMBER led on AI, yet both normalized distances
exceeded two. Both ensembles had $R_{\mathrm{NLL}}>1$ and
$R_{\mathrm{KL}}>1$ for every concept. Conversely, SNMF had low
$H_{\mathrm{test}}$ but a KL ratio near
or below one for every concept. Good target suppression with preserved
neighboring and SciQ accuracy therefore did not imply similarity to the twin
in these experiments. This disagreement is also expected from the definitions:
$H_{\mathrm{test}}$ compares accuracy with the full model and never uses the
twin, whereas the two ratios explicitly evaluate the twin as their reference.

As a robustness check against cancellation in the signed mean NLL difference,
we also measure question-level proximity to the twin. For concept $c$, let
\begin{equation}
\begin{aligned}
D_{\mathrm{abs}}(M,T;c)
&=\frac{1}{N}\sum_{i=1}^{N}|\ell_M(i)-\ell_T(i)|,\\
R_{\mathrm{abs}}(M,T;c)
&=\frac{D_{\mathrm{abs}}(M,T;c)}{D_{\mathrm{abs}}(F,T;c)}.
\end{aligned}
\end{equation}
Values of $R_{\mathrm{abs}}$ below one indicate that the erased model is closer
to the twin than the full model after taking the absolute difference on each
question. We additionally report $P_{\mathrm{closer}}$, the proportion of
held-out target questions for which the erased model is strictly closer to the
twin than the full model. Both statistics use the same 50 paired target-test
questions, with paired bootstrap confidence intervals.

\begin{table*}[t]
  \centering
  \scriptsize
  \begin{tabular}{llcc}
    \toprule
    Concept & Method & $R_{\mathrm{abs}}$ [95\% CI] $\downarrow$
      & $P_{\mathrm{closer}}$ [95\% CI] $\uparrow$ \\
    \midrule
    Ancient Rome & EMBER & 1.613 [1.151, 2.244] & 0.40 [0.26, 0.54] \\
                 & RMU   & \textbf{0.660 [0.513, 0.827]} & \textbf{0.64 [0.50, 0.76]} \\
                 & SNMF  & 0.940 [0.891, 0.986] & \textbf{0.64 [0.50, 0.76]} \\
                 & RMU+E & 1.384 [1.015, 1.873] & 0.44 [0.30, 0.58] \\
                 & SNMF+E & 1.558 [1.128, 2.143] & 0.42 [0.28, 0.56] \\
    \midrule
    Baseball     & EMBER & 2.128 [1.542, 2.978] & 0.26 [0.14, 0.38] \\
                 & RMU   & 1.227 [0.877, 1.680] & 0.46 [0.32, 0.60] \\
                 & SNMF  & \textbf{0.967 [0.888, 1.060]} & \textbf{0.66 [0.54, 0.78]} \\
                 & RMU+E & 2.079 [1.510, 2.903] & 0.34 [0.22, 0.48] \\
                 & SNMF+E & 2.241 [1.658, 3.101] & 0.28 [0.16, 0.40] \\
    \midrule
    AI           & EMBER & 2.112 [1.402, 3.167] & 0.38 [0.24, 0.52] \\
                 & RMU   & 1.006 [0.754, 1.344] & 0.54 [0.40, 0.68] \\
                 & SNMF  & \textbf{0.927 [0.876, 0.976]} & \textbf{0.64 [0.50, 0.76]} \\
                 & RMU+E & 2.196 [1.484, 3.251] & 0.34 [0.22, 0.48] \\
                 & SNMF+E & 2.152 [1.420, 3.239] & 0.40 [0.26, 0.54] \\
    \bottomrule
  \end{tabular}
  \caption{Question-level target NLL proximity to the concept-excluded twin.
  $R_{\mathrm{abs}}$ compares each method's mean absolute paired difference
  from the twin with the full model's; values below one indicate improvement.
  $P_{\mathrm{closer}}$ is the fraction of questions on which the method is
  closer than the full model. Intervals are paired-bootstrap percentile 95\%
  CIs over 50 questions (10,000 resamples, seed 0). Bold marks the best
  erasure condition in each concept and column.}
  \label{tab:question-level-similarity}
\end{table*}

The question-level results qualify the signed-mean ranking. RMU was clearly
closer than the full model for Rome ($R_{\mathrm{abs}}=0.660$), but its
intervals included one for Baseball and AI. In particular, RMU's AI signed
mean difference from the twin was only $-0.013$, whereas its absolute-distance
ratio was 1.006; positive and negative item-level differences had canceled.
SNMF had $R_{\mathrm{abs}}<1$ for all three concepts and was closer than the
full model on 64--66\% of target questions, although the Baseball ratio's
interval included one. EMBER was farther from the twin on average for every
concept, including ratios above two for Baseball and AI. Both ensembles were
also farther from the twin for every concept: all six $R_{\mathrm{abs}}$
intervals excluded one, and they were closer than the full model on only
28--44\% of questions. Thus, the
question-level analysis reinforces that high $H_{\mathrm{test}}$ does not
necessarily imply proximity to the twin.

\subsection{Selectivity and Variation Across Concepts}

To summarize whether changes were concentrated on the target, we computed the
post-hoc contrast between target NLL change from the full model and the average
change on neighboring and SciQ questions. EMBER's contrasts were 1.932 nats
per token for Rome, 1.242 for Baseball, and 2.124 for AI. The corresponding
values were 1.701, 1.126, and 2.325 for RMU+EMBER, and 1.886, 1.286, and 2.145
for SNMF+EMBER. All ensemble intervals excluded zero. RMU's contrasts were
$-0.094$, 0.151, and 0.011, while SNMF's were 0.131, 0.088, and 0.121.
EMBER and the ensembles therefore made large target-specific changes even
though they remained distant from the twins. The exact
definition, component contrasts, and uncertainty estimates are reported in
Appendix~\ref{app:diagnostics}.

Taken together, no method dominates under every interpretation of erasure.
EMBER and its ensembles produced the strongest target-specific changes, RMU
most closely matched
the twin's signed mean correct-answer NLL, and SNMF produced the lowest KL and
the most consistent question-level proximity to the twin. The change in
rankings across Rome, Baseball, and AI also shows that a
conclusion from one concept would not describe all three cases. Differences
in corpus coverage and overlap with neighboring topics are plausible
contributors, but the present measurements do not isolate their causal
effects.
